\subsection{Linear equations over a faithfully flat ring}

Let B be a ring and A a subring of B. We shall say that the ordered pair (A, B) has the linear extension property if it satisfies the following condition:

\begin{enumerate}
\def\labelenumi{(\Alph{enumi})}
\setcounter{enumi}{4}
\tightlist
\item
  Every solution \((y_{k})_{1\leqslant k\leqslant n}\) , consisting of elements of B, of a system of linear equations
\end{enumerate}

\[
\sum_ {k = 1} ^ {n} y _ {k} c _ {k \mathrm{l}} = d _ {i} \quad (1 \leqslant i \leqslant m)\tag{3}
\]

whose coefficients \(c_{kl}\) and right-hand sides \(d_{i}\) belong to A, is of the form

\[
y _ {k} = x _ {k} + \sum_ {j = 1} ^ {p} b _ {j} z _ {j k} \quad (1 \leqslant k \leqslant n)\tag{4}
\]

where \((x_{k})\) is a solution of (3) consisting of elements of A, the \(b_{j}\) belong to B and each of the \((z_{jk})_{1\leq k\leq n}\) is a solution of the homogeneous linear system associated with (3), consisting of elements of A.

PROPOSITION 13. Let A be a subring of a ring B. For the ordered pair (A, B) to satisfy the linear extension property, it is necessary and sufficient that B be a faithfully flat A-module.

The condition is sufficient. For, as B is a flat A-module, every solution with elements in B of the homogeneous linear system associated with (3) is a linear combination with coefficients in B of solutions consisting of elements of A ( \(\S\) 2, no. 11, Corollary 2, to Proposition 13). The problem then reduces to proving that the existence of a solution of (3) with elements in B implies the existence of one solution with elements in A. Now if we set

\[
c _ {k} = \left(c _ {k l}\right) _ {1 \leqslant i \leqslant m} \in \mathrm{A} _ {s} ^ {m}, \quad d = (d _ {i}) \in \mathrm{A} _ {s} ^ {m},
\]

system (3) is equivalent to the equation \(\sum_{k=1}^{n} y_k \otimes c_k = 1 \otimes d\) in \(\mathbf{B} \otimes_{\mathbf{A}} \mathbf{A}_s^m = \mathbf{B}_s^m\) . In other words, if \(\mathbf{M}\) is the sub-A-module of \(\mathbf{A}_s^m\) generated by the \(c_k\) ( \(1 \leqslant k \leqslant n\) ), the existence of the solution \((y_k)\) of (3) is equivalent (with the identifications made in no. 5) to the relation \(d \in \mathbf{BM} \cap \mathbf{A}_s^m\) ; but as \(\mathbf{BM} \cap \mathbf{A}_s^m = \mathbf{M}\) (no. 5, Proposition 10, (ii)), it implies \(d \in \mathbf{M}\) , that is, the existence of a solution \((x_k)\) of system (3) with elements in \(\mathbf{A}\) .

The condition is necessary. For suppose that (A, B) satisfies the linear extension property; we know already that B is a flat right A-module ( \(\S\) 2, no. 11, Corollary 2 to Proposition 13); we prove that, for every left ideal a of A, \(Ba \cap A = a\) , which shows that B is a faithfully flat right A-module (no. 5, Proposition 9, (d)). Now, let \(x \in Ba \cap A\) ; there exists by hypothesis \(y_i \in B\) and \(a_i \in a\) such that \(\sum_{i} y_i a_i = x\) ; property (E) applied to this linear equation with coefficients and right hand side in A shows that there exist \(x_i \in A\) such that \(x = \sum_{i} x_i a_i\) , hence \(x \in a\) .

